\section{分辨率折叠与稳定扇区：\texorpdfstring{$\varphi$--$\pi$--$e$}{varphi-pi-e} 三通道结构}

\subsection{稳定语言与 \texorpdfstring{$\varphi$}{varphi}-通道}

定义长度 $m$ 的黄金均值语言
\[
X_m=\{w\in\{0,1\}^m:\text{$w$ 不含子词 }11\}.
\]
计数满足 $|X_m|=F_{m+2}$，其指数增长率为 $\log\varphi$。该约束可视为“$\varphi$-通道”：它只依赖局部语法（禁止相邻 $1$），并将可允许状态空间的复杂度从 $2^m$ 压缩到 $F_{m+2}$。

在 Hilbert 表述中，取 $\mathcal{H}_m=\ell^2(\{0,1\}^m)$，定义正交投影
\[
(P_\varphi e_w)=
\begin{cases}
e_w,&w\in X_m,\\
0,&w\notin X_m,
\end{cases}
\]
则 $\mathrm{rank}(P_\varphi)=|X_m|$，从而得到稳定子空间的明确维数。

\subsection{分辨率折叠映射与退化谱}

有限观察者对 tick/地址的编码往往首先产生 $m$ 位原始串（可理解为“未规范化”表示）。将其压缩到稳定语言可通过分辨率折叠映射
\[
\Fold_m:\{0,1\}^m\to X_m,
\]
其原理可通过 Zeckendorf/Ostrowski 规范化实现：将原始串解释为 Fibonacci 权重表示并施加“无相邻 $1$”的规范化规则，从而得到唯一的规范形。该映射是满射且多对一，其平均退化度为
\[
d_m=\frac{2^m}{|X_m|}=\frac{2^m}{F_{m+2}},
\]
并呈 $(2/\varphi)^m$ 量级增长。

更强地，在特定 $m$ 可得到显式纤维分类。以 $m=6$ 为例，$2^6=64$，而 $|X_6|=F_8=21$，因此 $\Fold_6$ 将 64 个输入压缩为 21 个稳定词，并产生仅取 $\{2,3,4\}$ 的退化谱。该结构不仅提供压缩，也提供天然的容错/校验接口：不同原始输入在折叠后归并为同一稳定类型。

\subsection{\texorpdfstring{$\pi$}{pi}-通道：循环闭合与离散单值性}

仅有局部语法不足以保证“周期轨道兼容”。在符号动力系统中，长度 $m$ 词若要对应周期轨道片段，需要满足 wrap-around 的循环可容许性。对黄金均值约束而言，循环可容许性等价于排除 $w_mw_1=11$。定义
\[
X_m^{\mathrm{cyc}}=\{w\in X_m: w_mw_1\neq 11\},\qquad
X_m^{\mathrm{bdry}}=X_m\setminus X_m^{\mathrm{cyc}}.
\]
该分解体现“$\pi$-通道”：它以离散闭合测试替代连续情形的单值性/monodromy，并与周期轨道计数、$\zeta$ 恒等式相容。以 $m=6$ 为例出现 $18\oplus 3$ 的分裂：循环可容许者 18 个，边界态 3 个。该结构为后续的 $\zeta$ 正规化与稳定极点提供组合基础。

\subsection{\texorpdfstring{$e$}{e}-通道：指数正规化、\texorpdfstring{$\zeta$}{zeta} 极点与“极点屏障”}

在符号动力学中，周期轨道计数可由 Artin--Mazur $\zeta$ 函数编码。对黄金均值移位系统，其 $\zeta$ 函数具有
\[
\zeta(z)=\frac{1}{1-z-z^2}.
\]
分母因式分解给出极点 $z=\varphi^{-1}$ 与 $z=-\varphi$。取归一化 $z=r/\varphi$，则主极点被推至单位圆 $r=1$。该归一化在操作上构成“$e$-通道”：它以指数型尺度（几何级数/Abel 型核）将增长常数与可观测尺度对齐，从而在有限深度读出中形成可比较的稳定极点结构。该通道同时为噪声谱（如 $1/f$ 型）提供可计算的层级时间尺度阶梯，因为 Fibonacci/Zeckendorf 层级天然给出近几何的松弛时间族。

